%% =====================================================================
%%  B3-T-21-02 -- what B3 contributes to "The Surrender Tactic"
%%  -------------------------------------------------------------------
%%  LAYER A: APPEND-ONLY. Written once at the entry's write, then left
%%  alone. If the source has more to say later, add a bullet -- do not
%%  rewrite. Material found ONLY here goes in the `unique` slot with
%%  @unique set to yes, which makes it undeletable (rule R10).
%% =====================================================================
%% @kind: source
%% @id: B3-T-21-02
%% @book: B3
%% @topic: T-21-02
%% @locator: Law 22, pp. 165--174
%% @status: distilled
%% @unique: no
%% @integrated: yes
%% @updated: 2026-09-19
% @allow-rewrite: skeleton filled at the write of T-21-02 (planned -> stable lifecycle)

\dossier{B3}{T-21-02}{\bookshort{B3}}{Law 22, pp. 165--174}

\begin{distilled}
  \item When defeat is certain, surrender keeps you alive inside it: yield the field, conserve the core, let victory intoxicate the winner.
  \item The Brecht hearing: dressed as a clerk, agreeing the committee had a serious job, politely oblivious --- mocking the process from inside a pose of submission; released the same day while the posturing heroes were ruined for years.
  \item An opponent geared for resistance has planned against resistance; mildness makes those plans look absurd and gives them nothing to crush.
  \item The winner's vigilance rots on a schedule the survivor does not have to announce; the fight is relocated to a slower clock.
\end{distilled}
